\section[Modèles linéaires]{Modèles linéaires}

\begin{frame}{Approche théorique}
	Les modèles linéaires disposent d'un cadre \textbf{théorique} permettant de dériver ces intervalles ou ensembles. C'est ce que nous couvrirons dans cette section.
\end{frame}

\subsection[Régression linéaire]{Régression linéaire}

\begin{frame}{Hypothèses pour la régression linéaire}
	Pour la régression linéaire, les hypothèses déjà considérées sont encore d'actualité et même encore plus pertinentes:
	\begin{itemize}
		\item \textbf{Exogénéité}: $\E{\epsilon^{(i)} \mid X^{(i)}} = 0$
		\item \textbf{Homoscédasticité}: $\Var{\epsilon^{(i)} \mid X^{(i)}} = \sigeps$
		\item \textbf{Indépendance}: $\Cov{\epsilon^{(i)}, \epsilon^{(j)}} = 0 \quad \forall i \ne j$
	\end{itemize}
	Ces hypothèses sont appelées hypothèses de \textbf{Gauss-Markov}.
\end{frame}

\begin{frame}{Covariance de $\hat{\beta}$}
	Dans le cadre des hypothèses de Gauss-Markov (exogénéité, homoscédasticité et indépendance), la théorie de la régression des moindres carrés ordinaires, pour une application d'augmentation $\Phi: \mathcal{X} \to \tilde{\mathcal{X}}$, conduit à la matrice de covariance suivante pour $\hat{\beta}$ dans $\mathcal{M}_{d+1, d+1}\left(\mathbb{R}\right)$:
	\begin{equation}
		\Cov{\hat{\beta}} = \sigeps \left(\Phi(X)^T \Phi(X)\right)^{-1}
		\nonumber
	\end{equation}
	\pause
	En général, $\sigeps$ est inconnu. Il peut être estimé à partir de la variance des résidus $r = y - \hat{y}$:
	\begin{equation}
		s^2 = \frac{\text{SS}_{\text{res}}}{n - p} = \frac{\sum_{i=1}^n \left(y^{(i)} - \hat{y}^{(i)}\right)^2}{n - p}
		\nonumber
	\end{equation}
	Où $n$ est le nombre d'observations et $p = d + 1$ est le nombre de paramètres du modèle. $n - p$ est le nombre de degrés de liberté des résidus. Diviser par cette quantité permet d'obtenir un estimateur de variance $s^2$ non biaisé.
\end{frame}

\begin{frame}{Variance de $\hat{\beta}$}
	On peut déduire de la matrice de covariance la variance de $\hat{\beta}$ $\Var{\hat{\beta}} \in \mathcal{M}_{d+1, 1}\left(\mathbb{R^+}\right)$ en prenant les termes diagonaux de $\Cov{\hat{\beta}}$:
	\begin{equation}
		\Var{\hat{\beta}} = \diag \left(\Cov{\hat{\beta}}\right)
		\nonumber
	\end{equation}
	\pause
	L'erreur type $\ET[\hat{\beta}] \in \mathcal{M}_{d+1, 1}\left(\mathbb{R+}\right)$ peut alors se déduire facilement:
	\begin{equation}
		\ET[\hat{\beta}] = \sqrt{\Var{\hat{\beta}}}
		\nonumber
	\end{equation}
\end{frame}

\begin{frame}{Connexion à l'estimation mono-dimensionnelle}
	Dans la section précédente, nous avons vu que nous pouvions estimer l'erreur type à partir de l'écart type d'un échantillon:
	\begin{equation}
		\ET = \frac{\sigma}{\sqrt{n}}
		\nonumber
	\end{equation}
	Dans le modèle de régression linéaire, si nous prenons l'application d'augmentation $\Phi: x \mapsto 1$ alors $\beta$ est maintenant un scalaire:
	\begin{equation}
		\begin{aligned}
			\Var{\hat{\beta}} &= \sigeps \left(\Phi(X)^T \Phi(X)\right)^{-1} \\
			&= \sigeps \left(\begin{pmatrix}
				1 & 1 & \cdots & 1
			\end{pmatrix}\begin{pmatrix}
				1 \\ 1 \\ \vdots \\ 1
			\end{pmatrix}
			\right)^{-1} \\
			&= \sigeps \left(n\right)^{-1}
		\end{aligned}
		\nonumber
	\end{equation}
	Ce qui conduit à une erreur type semblable à celle vue auparavant $ET = \frac{\sigma_{\epsilon}}{\sqrt{n}}$	
\end{frame}

\begin{frame}{Variance de $\hat{Y}$}
	La variance de $\hat{Y} \mid X$, notée $\Var{\hat{Y} \mid X} \in \mathcal{M}_{n,1}(\mathbb{R})$ peut alors se calculer comme suit:
	\begin{equation}
		\Var{\hat{Y} \mid X} = \diag \left(\Phi(X) \Cov{\hat{\beta}} \Phi(X)^T\right)
		\nonumber
	\end{equation}
	\pause
	Et par conséquent, on a l'erreur type suivante:
	\begin{equation}
		\ET[\hat{Y} \mid X] = \sqrt{\Var{\hat{Y} \mid X}}
		\nonumber
	\end{equation}	
	\pause
	En supposant que les $\epsilon^{(i)}$ suivent une loi normale $\mathcal{N}(0, \sigeps)$ on peut alors déterminer un intervalle de confiance, soit sur $\hat{\beta}$, soit sur $\hat{Y}$.
\end{frame}

\begin{frame}{Intervalles de confiance}
	Notons $\varphi$ la fonction de répartition d'une loi normale standard $\mathcal{N}(0, 1)$. L'intervalle de confiance d'un niveau $1-\alpha$, noté $\text{IC}_{1-\alpha}$ pour $\hat{\beta}$ est donc:
	\begin{equation}
		\text{IC}_{1-\alpha}\left(\hat{\beta}\right) = \left[\hat{\beta} \pm \varphi^{-1}(1-\alpha / 2) \sqrt{\Var{\hat{\beta}}}\right]
		\nonumber
	\end{equation}
	\pause
	Pour la prédiction $\hat{Y} \mid X$, l'intervalle de confiance d'un niveau $1-\alpha$ est:
	\begin{equation}
		\begin{split}
			\text{IC}_{1-\alpha}\left(\hat{Y} \mid X\right) = \Biggl[ \hat{Y} \pm \varphi^{-1}(1-\alpha / 2) \sqrt{\Var{\hat{Y} \mid X}} \Biggr]
		\end{split}
		\nonumber
	\end{equation}
	\pause
	Note: Comme $\sigeps$ est généralement inconnue, il est classique d'utiliser une loi de Student avec $n - p$ degrés de liberté $\mathcal{T}(n - p)$ à la place de la loi normale standard.
\end{frame}

\begin{frame}{Intervalle de prédiction}
	L'intervalle de prédiction d'un niveau $1-\alpha$, noté $\text{IP}_{1-\alpha}$, pour la prédiction $\hat{Y} \mid X$ est:
	\begin{equation}
		\begin{split}
			\text{IP}_{1-\alpha}\left(\hat{Y} \mid X\right) = \Biggl[ \hat{Y} \pm \varphi^{-1}(1-\alpha / 2) \sqrt{\sigeps + \Var{\hat{Y} \mid X}} \Biggr]
		\end{split}
		\nonumber
	\end{equation}
	\pause
	Noter qu'il suffit d'ajouter $\sigeps$ à la variance de $\hat{Y} \mid X$ pour passer de l'intervalle de confiance à l'intervalle de prédiction.
\end{frame}

\subsection[Régression logistique]{Régression logistique}

\begin{frame}{Définitions supplémentaires}
	Nous aurons besoin de quelques définitions complémentaires associées à des problèmes de classification.
	\begin{block}{Définition: cote}
		La \textbf{cote} est un nombre réel dans $[0, +\infty[$ qui représente le rapport entre une probabilité d'un événement et la probabilité complémentaire de ce même événement:
		\begin{equation}
			\cote = \frac{p}{1 - p}
			\nonumber
		\end{equation}
	\end{block}
	\pause
	\begin{itemize}
		\item Si $p=0.5$, la cote est de $1:1$ (une chance sur deux).
		\item Si $p=0.8$, la cote est de $4:1$ (quatre contre un).
	\end{itemize}	
\end{frame}

\begin{frame}{Définitions supplémentaires}
	\begin{block}{Définition: logit}
		Le \textbf{logit} est un nombre réel dans $]-\infty, +\infty[$ qui représente le logarithme d'une cote:
		\begin{equation}
			\logit{p} = \ln{\left(\frac{p}{1 - p}\right)}
			\nonumber
		\end{equation}
	\end{block}
	\pause
	\begin{itemize}
		\item Si $p=0.5$, la cote est de $1:1$ et le logit est de $0$.
		\item Si $p=0.8$, la cote est de $4:1$ et le logit est approximativement $1.387$.
	\end{itemize}	
\end{frame}

\begin{frame}{Définitions supplémentaires}
	\begin{block}{Définition: score prédit}
		Dans un problème de régression logistique, on définit le \textbf{score}, un nombre réel dans $]-\infty, +\infty[$, comme la prédiction réalisée par le modèle linéaire:
		\begin{equation}
			\score{x} = \Phi(x) \hat{\beta}
			\nonumber
		\end{equation}
	\end{block}
\end{frame}

\begin{frame}{Définitions supplémentaires}
	\begin{block}{Définition: probabilité prédite}
		Dans un problème de régression logistique, on définit la \textbf{probabilité}, notée $\hat{p} \in ]0, 1[$, comme la prédiction réalisée par le modèle linéaire transformée par la fonction logistique $\sigma$:
		\begin{equation}
			\hat{p}(x) = \sigma\left(\score{x}\right) = \sigma\left(\Phi(x) \hat{\beta}\right)
			\nonumber
		\end{equation}
	\end{block}
\end{frame}

\begin{frame}{Dérivation}
	Calculer le logit de la probabilité prédite revient à effectuer le calcul suivant:
	\begin{equation}
		\begin{aligned}
			\logit{\hat{p}(x)}
			&= \ln{\left(\frac{\hat{p}(x)}{1 - \hat{p}(x)}\right)}
			= \ln{\left(\frac{\sigma(\score{x})}{1 - \sigma(\score{x})}\right)} \\
			&= \ln{\left(\frac{1}{1 + \exp{(-\score{x})}} \frac{1}{1 - \frac{1}{1 + \exp{(-\score{x})}}}\right)} \\
			&= \ln{\left(\frac{1}{1 + \exp{(-\score{x})}} \frac{1}{\frac{\exp{(-\score{x})}}{1 + \exp{(-\score{x})}}}\right)} \\
			&= \ln{\left(\frac{1}{\exp{(-\score{x})}}\right)} = \score{x}
		\end{aligned}
		\nonumber
	\end{equation}
\end{frame}

\begin{frame}{Dérivation}
	Dans un problème de régression logistique, le score et le logit représentent la \textit{même} quantité. Cela traduit la nature linéaire de la régression logistique: c'est une régression linéaire sur le score, \textit{enrobée} dans une fonction logistique pour calculer des probabilités !
\end{frame}

\begin{frame}{Revisite de l'incertitude en régression linéaire}
	En régression linéaire, nous avons pris l'hypothèse que la variance $\sigeps$ est constante (c'est notre hypothèse d'homoscédasticité). Nous avons écrit que la matrice de covariance est:
	\begin{equation}
		\Cov{\hat{\beta}} = \sigeps \left(\Phi(X)^T \Phi(X)\right)^{-1}
		\nonumber
	\end{equation}
	En introduisant une matrice diagonale $W \in \mathcal{M}_{n,n}(\mathbb{R})$:
	\begin{equation}
		W = \frac{1}{\sigeps} I_n
		\nonumber
	\end{equation}
	On peut dériver une formulation équivalente de la covariance des estimations $\beta$:
	\begin{equation}
		\Cov{\hat{\beta}} = \sigeps \left(\Phi(X)^T \Phi(X)\right)^{-1} = \left(\Phi(X)^T W \Phi(X)\right)^{-1}
		\nonumber
	\end{equation}	
\end{frame}

\begin{frame}{Incertitude en régression logistique}
	\begin{itemize}
		\item Comme la régression logistique résout un problème de régression linéaire sur le logit (ou le score), il est naturel de définir une matrice de poids $W$ basée sur la variance du logit.
		\pause
		\item On rappelle que chaque observation est associée à une loi de Bernoulli. Mais comme chaque paramètre de la loi de Bernoulli est différent, on doit utiliser une variance différente pour chaque observation.
		\pause
		\item La variance d'une variable aléatoire $Y^{(i)}$ suivant une loi de Bernoulli de paramètre $p^{(i)}$ est donnée par $p^{(i)} \left(1 - p^{(i)}\right)$.
		\pause
		\item On pourrait montrer que la variance de $\logit{p^{(i)}}$ est alors: 
		\begin{equation}
			\Var{\logit{p^{(i)}}} = \frac{1}{p^{(i)}\left(1 - p^{(i)}\right)}
			\nonumber
		\end{equation}
	\end{itemize}
\end{frame}

\begin{frame}{Incertitude en régression logistique}
	On peut donc définir, pour la régression logistique, une nouvelle matrice de poids $W$ telle que:
	\begin{equation}
		W_{ij} = \begin{cases}
			0 & \text{si} \, i \ne j \\
			\hat{p}^{(i)}\left(1 - \hat{p}^{(i)}\right) & \text{sinon}
		\end{cases}
		\nonumber
	\end{equation}
	\pause
	Et on en déduit la matrice de covariance de $\hat{\beta}$ pour la régression logistique:
	\begin{equation}
		\Cov{\hat{\beta}} = \left(\Phi(X)^T W \Phi(X)\right)^{-1}
		\nonumber
	\end{equation}
	\pause
	Et la variance de $\hat{\beta}$:
	\begin{equation}
		\Var{\hat{\beta}} = \diag \left(\Cov{\hat{\beta}}\right)
		\nonumber
	\end{equation}	
\end{frame}

\begin{frame}{Intervalle de confiance sur $\hat{\beta}$}
	Comme il s'agit de la partie linéaire du modèle pour calculer le score, on peut faire la même hypothèse que pour le modèle de régression linéaire, c'est à dire supposer que les incertitudes sur le score sont de type Gaussien. L'intervalle de confiance au niveau $1 - \alpha$ sur $\hat{\beta}$ est alors:
	\begin{equation}
		\text{IC}_{1-\alpha}\left(\hat{\beta}\right) = \left[\hat{\beta} \pm \varphi^{-1}(1 - \alpha/2) \sqrt{\Var{\hat{\beta}}}\right]
		\nonumber
	\end{equation}	
\end{frame}

\begin{frame}{Intervalle de confiance sur le logit ou le score}
	Pour une nouvelle observation $x^{(n+1)}$, on obtient d'abord la variance sur le score linéaire (le logit):
	\begin{equation}
		\Var{\logit{\hat{p}\left(x^{(n+1)}\right)}} = \diag\left(\Phi(x^{(n+1)}) \Cov{\hat{\beta}} \Phi(x^{(n+1)})^T\right)
		\nonumber
	\end{equation}
	\pause
	Puis, l'intervalle de confiance au niveau $1-\alpha$ sur le score (ou logit) se définit donc comme:
	\begin{equation}
		\begin{split}
			\text{IC}_{1-\alpha}\left(\logit{\hat{p}\left(x^{(n+1)}\right)}\right) = \bigg[
			&\logit{\hat{p}\left(x^{(n+1)}\right)} \\
			&\pm \varphi^{-1}(1 - \alpha/2) \sqrt{\Var{\logit{\hat{p}\left(x^{(n+1)}\right)}}}\bigg]			
		\end{split}
		\nonumber
	\end{equation}
\end{frame}

\begin{frame}{Intervalle de confiance sur la probabilité}
	En appliquant la fonction logistique $\sigma$ à l'intervalle précédent on obtient l'intervalle de confiance au niveau $1 - \alpha$ de la probabilité $\hat{p}(x^{(n+1)})$:
	\begin{equation}
		\begin{aligned}
			&\hat{p}_{\text{inf}} = \sigma\left(\logit{\hat{p}\left(x^{(n+1)}\right)} - \varphi^{-1}(1 - \alpha/2) \sqrt{\Var{\logit{\hat{p}\left(x^{(n+1)}\right)}}}\right) \\
			&\hat{p}_{\text{sup}} = \sigma\left(\logit{\hat{p}\left(x^{(n+1)}\right)} + \varphi^{-1}(1 - \alpha/2) \sqrt{\Var{\logit{\hat{p}\left(x^{(n+1)}\right)}}}\right) \\
			&\text{IC}_{1-\alpha}\left(\hat{p}\left(x^{(n+1)}\right)\right) = \left[\hat{p}_{\text{inf}}, \hat{p}_{\text{sup}}\right]
		\end{aligned}
		\nonumber
	\end{equation}
\end{frame}

\begin{frame}{Région de prédiction sur la classe}
	Dans un problème de classification binaire, $\hat{p}\left(x^{(n+1)}\right)$ représente en fait $P(y^{(n+1)} = 1 \mid x^{(n+1)})$. \\
	En fonction de la valeur de $\hat{p}_{\text{inf}}$ et de $\hat{p}_{\text{sup}}$ on peut alors définir la région de prédiction au niveau $1 - \alpha$ comme suit:
	\begin{equation}
		C\left(x^{(n+1)}\right) = 
		\begin{cases}
			\{0\} & \text{si} \, \hat{p}_{\text{sup}} \le \alpha \\
			\{1\} & \text{si} \, \hat{p}_{\text{inf}} \ge 1 - \alpha \\
			\{0, 1\} & \text{sinon}
		\end{cases}
		\nonumber
	\end{equation}
	\pause
	\begin{itemize}
		\item Si $\text{IC}_{95\%}\left(\hat{p}\left(x^{(n+1)}\right)\right) = [0.92, 0.98]$ alors $C\left(x^{(n+1)}\right) = \{0, 1\}$.
		\item Si $\text{IC}_{95\%}\left(\hat{p}\left(x^{(n+1)}\right)\right) = [0.01, 0.03]$ alors $C\left(x^{(n+1)}\right) = \{0\}$.
		\item Si $\text{IC}_{95\%}\left(\hat{p}\left(x^{(n+1)}\right)\right) = [0.98, 0.99]$ alors $C\left(x^{(n+1)}\right) = \{1\}$.
	\end{itemize}
\end{frame}
